\documentclass[blue]{beamer}
\mode<presentation> {
  \setbeamertemplate{background canvas}[vertical shading][bottom=red!10,top=blue!10]
  \usetheme{Boadilla}
  \usefonttheme[onlysmall]{serif}
}
\usepackage{pgf,pgfarrows,pgfnodes,pgfautomata,pgfheaps,pgfshade}
\usepackage{amsmath,amssymb,bm}
\usepackage{colortbl}
\usepackage[english]{babel}
\usepackage[utf8]{inputenc}
\setbeamercovered{dynamic}
\def\R{\mathbb{R}}
\def\Q{\mathbb{Q}}
\def\N{\mathbb{N}}
\def\E{\mathbb{E}}
\newcommand{\ps}[2]{ \langle #1 , #2 \rangle }
\def\1{\mathbf{1}}
\newcommand{\nor}[1]{\left\| #1 \right\|}
\def\leq{\leqslant}
\def\geq{\geqslant}
\def\et{\quad \textrm{and} \quad}
\def\ou{\quad \textrm{or} \quad}
\def\boite{\hskip 0.5mm {\Box} \hskip 0.5mm}
\DeclareMathOperator\PPE{PPE}
\DeclareMathOperator\vNM{vNM}
\DeclareMathOperator\argmax{argmax}
\DeclareMathOperator\Prob{Prob}
\DeclareMathOperator\supp{supp}
\DeclareMathOperator\Eq{Eq}
\DeclareMathOperator\inte{int}
\DeclareMathOperator\co{co}
\DeclareMathOperator\cl{cl}
\DeclareMathOperator\Dirac{Dirac}
\DeclareMathOperator\F{F}
\DeclareMathOperator\As{As}
\DeclareMathOperator\Ir{Ir}
\DeclareMathOperator\MRS{MRS}
\DeclareMathOperator\Range{Im}
\DeclareMathOperator\DIV{DIV}
%------------------------------ theorem styles
\theoremstyle{definition}
\newtheorem*{assumption}{Assumption}
\newtheorem*{conjecture}{Conjecture}
\theoremstyle{example}
\newtheorem*{proposition}{Proposition}
\newtheorem*{theoremns}{Theorem}
\newtheorem*{remark}{Remark}
%-----------------------------------------
\title[Corporate Finance]{Theory of Corporate Finance}
\author[Victor Filipe]{V. Filipe Martins-da-Rocha}
\institute[FGV EESP]{Escola de Economia de S\~ao Paulo}
\date[April--June, 2026]{Part 1.a. Stock Market Equilibrium: Objective of the Firm}
\subject{Economic Theory}
\begin{document}
\frame{\titlepage}

%--------------------------------------------------------------------
\begin{frame}\frametitle{Outline}
\begin{itemize}
\item A two-period economy: investors and a firm
\item Stock-market equilibrium
\item Asset pricing at equilibrium
    \begin{itemize}
    \item Fundamental asset-pricing theorem
    \item Personalized state prices
    \end{itemize}
\item The objective of the firm
    \begin{itemize}
    \item Separation of ownership and control
    \item Competitive price perceptions
    \end{itemize}
\item Unanimity and value maximization
    \begin{itemize}
    \item Unanimity under complete markets
    \item Spanning condition
    \item Strong unanimity
    \end{itemize}
\end{itemize}
\end{frame}

%====================================================================
\section{The Model}
%====================================================================

\begin{frame}\frametitle{A simple intertemporal model}
\begin{block}{Time}
Two dates: $t \in \{0, 1\}$
\end{block}
\begin{block}{Uncertainty at $t=1$}
A finite set $S$ of states of nature
\end{block}
\begin{block}{Commodity}
A single divisible commodity, available at every date and state
\end{block}
\end{frame}

%--------------------------------------------------------------------
\begin{frame}\frametitle{Agents}
\begin{block}{Investors}
A finite set $I$ of risk-averse, expected-utility investors
\end{block}
\begin{block}{Firm}
One firm, privately owned by the investors
\end{block}
\begin{block}{Financial assets}
A finite set $A$ of state-contingent claims (securities)
\end{block}
\end{frame}

%--------------------------------------------------------------------
\begin{frame}\frametitle{The firm's technology}
\begin{block}{Production set}
\begin{equation*}
Y \subset \R \times \R^{S}
\end{equation*}
\end{block}
\begin{block}{Production plan}
A vector $y = (y_0, y_1) \in Y$ with $y_1 = (y_s)_{s \in S}$
\end{block}
\end{frame}

%--------------------------------------------------------------------
\begin{frame}\frametitle{No instantaneous production}
\begin{assumption}
The firm invests at $t=0$ and produces at $t=1$:
\begin{equation*}
Y \subset \R_{-} \times \R^{S}_{+}
\end{equation*}
\end{assumption}
\bigskip

For $y \in Y$:
\begin{itemize}
\item $-y_0 \geq 0$: input at $t=0$ (investment)
\item $y_s \geq 0$: output at $t=1$ in state $s$
\end{itemize}
\end{frame}

%--------------------------------------------------------------------
\begin{frame}\frametitle{Shares and ownership}
The firm is privately owned: investor $i$ holds initial share
\begin{equation*}
\xi^i \geq 0,
\qquad
\sum_{i \in I} \xi^i = 1
\end{equation*}

\begin{itemize}
\item $\xi^i = 0$: investor $i$ is not initially a shareholder
\item $\xi^i = 1$: investor $i$ is the sole owner
\end{itemize}
\end{frame}

%--------------------------------------------------------------------
\begin{frame}\frametitle{Stock market}
At $t=0$, agents may trade ownership through a competitive stock market
\bigskip

\begin{itemize}
\item The firm's equity trades at a linear price $E$ (in units of the commodity)
\item Investors take $E$ as given
\item Buying a fraction $\alpha \geq 0$ costs $E\alpha$
\end{itemize}
\end{frame}

%--------------------------------------------------------------------
\begin{frame}\frametitle{Financial assets}
Each asset $a \in A$ is a state-contingent claim
\begin{equation*}
\kappa(a) = (\kappa_s(a))_{s \in S} \in \R^S
\end{equation*}
\bigskip

A \emph{portfolio} is $\theta \in \R^A$:
\begin{itemize}
\item $\theta(a) \geq 0$: long position
\item $\theta(a) \leq 0$: short position
\end{itemize}
\bigskip

At $t=1$ in state $s$, the holder receives $\sum_a \kappa_s(a)\theta(a)$
\bigskip

An asset $a$ is \emph{risk-free} if $\kappa_s(a) = 1$ for all $s \in S$
\end{frame}

%--------------------------------------------------------------------
\begin{frame}\frametitle{Competitive security markets}
At $t=0$ each security trades at a linear price $q(a)$
\bigskip

Buying portfolio $\theta$ costs
\begin{equation*}
q \cdot \theta = \sum_{a \in A} q(a)\theta(a)
\end{equation*}
\bigskip

If investor spends $w(a) = q(a)\theta(a)$ in asset $a$, the gross return at state $s$ is
\begin{equation*}
R_s(a) = \frac{\kappa_s(a)}{q(a)},
\qquad
r_s(a) = R_s(a) - 1
\end{equation*}
\end{frame}

%--------------------------------------------------------------------
\begin{frame}\frametitle{Investors: characteristics}
For each investor $i \in I$:
\begin{itemize}
\item initial wealth $\omega^i_0 \geq 0$ at $t=0$
\item state-contingent endowment $\omega^i_1 \in \R^S_+$
\item initial share $\xi^i \geq 0$
\item Bernoulli utility $u^i : \R_+ \to [-\infty, \infty)$
\item discount factor $\beta^i > 0$
\item subjective beliefs $\nu^i \in \Prob(S)$
\end{itemize}
\end{frame}

%--------------------------------------------------------------------
\begin{frame}\frametitle{Investors: preferences}
At $t=0$, investor $i$ chooses $(c_0, c_1) \in \R_+ \times \R^S_+$
\bigskip

\begin{block}{Expected utility}
\begin{equation*}
U^i(c_0, c_1) = u^i(c_0) + \beta^i \sum_{s \in S} \nu^i_s u^i(c_s)
\end{equation*}
\end{block}
\end{frame}

%--------------------------------------------------------------------
\begin{frame}\frametitle{The firm's actions}
At $t=0$, taking security prices $q$ as given, the firm chooses
\bigskip

\begin{itemize}
\item a production plan $y = (y_0, y_1) \in Y$
\item a financial plan $\beta \in \R^A$:
    \begin{itemize}
    \item $\beta(a) > 0$: firm \emph{issues} $\beta(a)$ units of asset $a$
    \item $\beta(a) < 0$: firm \emph{purchases} $|\beta(a)|$ units of asset $a$
    \end{itemize}
\end{itemize}
\end{frame}

%--------------------------------------------------------------------
\begin{frame}\frametitle{The cash-flow stream}
The plan $(y, \beta)$ generates a cash-flow stream $\delta = (\delta_0, \delta_1)$:
\bigskip

\begin{block}{At $t=0$}
\begin{equation*}
\delta_0 = y_0 + q \cdot \beta
\end{equation*}
Owners share $\delta_0$ in proportion to $\xi^i$
\end{block}

\begin{block}{At $t=1$ in state $s$}
\begin{equation*}
\delta_s = y_s - \kappa_s \cdot \beta
\end{equation*}
Shareholders at $t=1$ (with $\alpha^i > 0$) share $\delta_s$
\end{block}
\end{frame}

%--------------------------------------------------------------------
\begin{frame}\frametitle{Ownership structures: I}
\begin{block}{Sole proprietorship}
\begin{itemize}
\item No stock market for the firm
\item One agent $i^\ast$ owns the entire firm: $\xi^{i^\ast} = 1$
\item The owner manages in his own interest
\end{itemize}
\end{block}

\begin{block}{Partnership}
\begin{itemize}
\item No stock market for the firm
\item Several shareholders
\item Possible disagreement on the firm's objective
\end{itemize}
\end{block}
\end{frame}

%--------------------------------------------------------------------
\begin{frame}\frametitle{Ownership structures: II}
\begin{block}{Corporation}
\begin{itemize}
\item Active stock market for the firm's equity
\item \textcolor{red}{Separation of ownership and control}
\item Decisions delegated to a manager whose mandate is to be defined
\end{itemize}
\end{block}
\bigskip

This is the structure we focus on
\end{frame}

%--------------------------------------------------------------------
\begin{frame}\frametitle{Limited liability}
Suppose the firm finances itself only through a risk-free bond
\bigskip

\begin{itemize}
\item At $t=0$, the firm issues debt by selling $Z$ units of the bond
\item The firm is then said to be \emph{levered}
\item Shareholders are protected by \emph{limited liability}: they never have to deliver
\item In states $s$ where $\delta_s < 0$, a \emph{bankruptcy rule} applies and the entire output $y_s$ is distributed to creditors
\end{itemize}
\bigskip

This chapter abstracts from default; bankruptcy is studied in subsequent parts
\end{frame}

%====================================================================
\section{Stock-Market Equilibrium}
%====================================================================

\begin{frame}\frametitle{Investor's budget set: setup}
Each investor $i$ takes as given (and perfectly anticipates):
\begin{itemize}
\item the firm's production--finance plan $(y, \beta)$
\item the security and equity prices $(q, E)$
\end{itemize}
\bigskip

He chooses $(c, \theta, \alpha) \in \R_+ \times \R^S_+ \times \R^A \times \R_+$
\end{frame}

%--------------------------------------------------------------------
\begin{frame}\frametitle{Investor's budget set: constraints}
\begin{block}{At $t=0$}
\begin{equation*}
c_0 + q \cdot \theta + E\alpha \leq \omega^i_0 + (E + \delta_0)\xi^i
\end{equation*}
\end{block}

\begin{block}{At $t=1$ in state $s$}
\begin{equation*}
c_s \leq \omega^i_s + \kappa_s \cdot \theta + \delta_s \alpha
\end{equation*}
\end{block}
\bigskip

The set of admissible actions is denoted $B^i(q, E)$
\end{frame}

%--------------------------------------------------------------------
\begin{frame}\frametitle{Stock-market equilibrium}
\begin{block}{Definition}
A stock-market equilibrium is a list
\begin{equation*}
\bigl\{ (y, \beta), (q, E), (c^i, \theta^i, \alpha^i)_{i \in I} \bigr\}
\end{equation*}
of:
\begin{itemize}
\item production--finance plan $(y, \beta) \in Y \times \R^A$
\item prices $(q, E) \in \R^A \times \R$
\item investor actions $(c^i, \theta^i, \alpha^i)$ for each $i \in I$
\end{itemize}
\end{block}
\bigskip

such that the conditions on the next slide are satisfied
\end{frame}

%--------------------------------------------------------------------
\begin{frame}\frametitle{Stock-market equilibrium: conditions}
\begin{itemize}
\item[(a)] each investor optimizes:
\begin{equation*}
(c^i, \theta^i, \alpha^i) \in \argmax \bigl\{ U^i(c) \colon (c, \theta, \alpha) \in B^i(q, E) \bigr\}
\end{equation*}

\item[(b)] securities and stock markets clear:
\begin{equation*}
\sum_{i \in I} \alpha^i = 1 \et \beta(a) = \sum_{i \in I} \theta^i(a), \quad \text{for all $a\in A$}
\end{equation*}

\item[(c)] commodity markets clear:
\begin{equation*}
\sum_{i \in I} c^i = \sum_{i \in I} \omega^i + y
\end{equation*}
\end{itemize}
\end{frame}

%--------------------------------------------------------------------
\begin{frame}\frametitle{Standard assumptions}
For each investor $i \in I$:
\begin{itemize}
\item[(C.1)] $u^i$ continuously differentiable, strictly increasing, strictly concave
\item[(C.2)] endowments are strictly positive: $\omega^i_0 > 0$ and $\omega^i_s > 0$, for all $s\in S$
\item[(C.3)] every state has positive probability: $\nu^i_s > 0$, for all $s\in S$
\item[(C.4)] individually rational consumption plans are interior:
\begin{equation*}
U^i(c) \geq U^i(\omega^i) \Longrightarrow c_0 > 0 \et c_s > 0, \quad \text{for all $s\in S$}
\end{equation*}
\end{itemize}
\end{frame}

%====================================================================
\section{Asset Pricing}
%====================================================================

\begin{frame}\frametitle{Marginal rate of substitution}
For each investor $i$ with interior optimal consumption $c^i$:
\bigskip

\begin{block}{Marginal rate of substitution}
\begin{equation*}
\MRS^i_s(c^i) := \beta^i \nu^i_s \frac{(u^i)'(c^i_s)}{(u^i)'(c^i_0)}
\end{equation*}
\end{block}
\bigskip

Under common beliefs ($\nu^i = \nu$), the quantity
\begin{equation*}
\beta^i \frac{(u^i)'(c^i_s)}{(u^i)'(c^i_0)}
\end{equation*}
is investor $i$'s \emph{stochastic discount factor} relative to $\nu$
\end{frame}

%--------------------------------------------------------------------
\begin{frame}\frametitle{Fundamental asset-pricing theorem}
\begin{theoremns}
Under (C.1)--(C.3), at any stock-market equilibrium there exists a vector of \emph{state-price deflators}
\begin{equation*}
\lambda = (\lambda_s)_{s \in S} \in \R^{S}_{++}
\end{equation*}
such that
\begin{equation*}
q(a) = \sum_{s \in S} \lambda_s \kappa_s(a), \quad \text{for all $a \in A$}
\end{equation*}
and
\begin{equation*}
E = \sum_{s \in S} \lambda_s \delta_s
\end{equation*}
\end{theoremns}
\bigskip

This is a no-arbitrage statement: positive equilibrium prices admit a strictly positive linear pricing functional
\end{frame}

%--------------------------------------------------------------------
\begin{frame}\frametitle{Personalized state prices}
\begin{proposition}
Under (C.1)--(C.4), at any stock-market equilibrium and for every investor $i \in I$:
\begin{equation*}
q(a) = \sum_{s \in S} \MRS^i_s \kappa_s(a), \quad \text{for all $a \in A$}
\end{equation*}
For every shareholder of the firm at $t=1$ ($\alpha^i > 0$):
\begin{equation*}
E = \sum_{s \in S} \MRS^i_s \delta_s
\end{equation*}
\end{proposition}
\bigskip

The proof uses the first-order conditions of investor $i$'s problem, which hold with equality under (C.4)
\end{frame}

%--------------------------------------------------------------------
\begin{frame}\frametitle{Uniqueness of $\lambda$}
The state-price deflator $\lambda$ is generally \emph{not unique}
\bigskip

\begin{block}{Incomplete markets}
If $\Range \kappa \subsetneq \R^S$, the no-arbitrage equation admits multiple positive solutions, and different investors may have different $\MRS^i$
\end{block}

\begin{block}{Complete markets}
If $\Range \kappa = \R^S$, the equation has a \emph{unique} solution, and the personalized-pricing proposition forces
\begin{equation*}
\MRS^i = \lambda \quad \text{for every $i\in I$}
\end{equation*}
\end{block}
\end{frame}

%====================================================================
\section{The Objective of the Firm}
%====================================================================

\begin{frame}\frametitle{The objective of the firm}
\begin{center}
\Large \alert{When choosing $(y, \beta)$, what should the firm maximize?}
\end{center}
\bigskip
\bigskip

The answer must come from \emph{shareholders' preferences}
\end{frame}

%--------------------------------------------------------------------
\begin{frame}\frametitle{An arbitrary candidate}
A natural candidate: maximize the firm's value
\begin{equation*}
V = E + \delta_0 = E + q \cdot \beta + y_0
\end{equation*}
\bigskip

This is the equity price plus the date-0 cash flow
\bigskip

But why this particular objective rather than, say, a weighted sum of shareholders' utilities?
\end{frame}

%--------------------------------------------------------------------
\begin{frame}\frametitle{Separation of ownership and control}
\begin{itemize}
\item The firm belongs to its shareholders $\{i \colon \xi^i > 0\}$
\item Each shareholder $i$ would prescribe his preferred plan $(y^i, \beta^i)$
\item Different shareholders may want \emph{different} plans
\bigskip
\item If all shareholders agree on a common objective, we have \emph{unanimity}
\item Under unanimity, we can \textcolor{red}{separate ownership from control}
\end{itemize}
\bigskip

\textbf{Question:} Under what conditions does unanimity arise?
\end{frame}

%--------------------------------------------------------------------
\begin{frame}\frametitle{Impact of a perturbation}
Fix an equilibrium and consider a small change $(\Delta y, \Delta \beta)$
\bigskip

Investor $i$'s presumed change in utility:
\begin{equation*}
\frac{\Delta^i U^i}{(u^i)'(c^i_0)} \approx \xi^i [\Delta^i E + \Delta \delta_0]
+ \alpha^i \Bigl[ \sum_{s \in S} \MRS^i_s \Delta \delta_s - \Delta^i E \Bigr]
\end{equation*}
\bigskip

where $\Delta^i E$ is investor $i$'s \emph{presumed} effect on the equity price
\end{frame}

%--------------------------------------------------------------------
\begin{frame}\frametitle{Competitive price perceptions}
\begin{block}{Definition (Grossman--Hart 1979)}
Each investor $i$ values any state-contingent income stream $w \in \R^S$ at
\begin{equation*}
\sum_{s \in S} \MRS^i_s w_s
\end{equation*}
using his own marginal rates of substitution as state-price deflators
\end{block}
\bigskip

\textbf{Interpretation:} given $c^i$, investor $i$ is indifferent between buying or selling a small claim paying $w$ at this price
\end{frame}

%--------------------------------------------------------------------
\begin{frame}\frametitle{Market consistency}
If $w$ is marketed, $w_s = \kappa_s \cdot \theta_w$, then
\begin{equation*}
\sum_{s \in S} \MRS^i_s w_s = q \cdot \theta_w
\end{equation*}
\bigskip

\textbf{All investors agree} on the value of marketed income streams, even when their MRSs differ
\bigskip

The disagreement only concerns \emph{non-marketed} state-contingent streams
\end{frame}

%--------------------------------------------------------------------
\begin{frame}\frametitle{Impact under competitive price perceptions}
Under competitive price perceptions, investor $i$ presumes
\begin{equation*}
\Delta^i E = \sum_{s \in S} \MRS^i_s \Delta \delta_s
\end{equation*}
\bigskip

Substituting into the impact formula:
\begin{equation*}
\Delta^i U^i \approx (u^i)'(c^i_0) \cdot \xi^i \cdot \Bigl[ \Delta \delta_0 + \sum_{s \in S} \MRS^i_s \Delta \delta_s \Bigr]
\end{equation*}
\bigskip

\begin{block}{Present-value criterion}
A shareholder favors any change that maximizes the present value of the firm's dividend stream evaluated at his own MRS
\end{block}
\end{frame}

%--------------------------------------------------------------------
\begin{frame}\frametitle{Disagreement under incomplete markets}
If $\MRS^{i_1} \neq \MRS^{i_2}$, a perturbation $(\Delta y, \Delta \beta)$ may yield
\begin{equation*}
\sum_{s} \MRS^{i_1}_s \Delta \delta_s + \Delta \delta_0 > 0
\end{equation*}
\begin{equation*}
\sum_{s} \MRS^{i_2}_s \Delta \delta_s + \Delta \delta_0 < 0
\end{equation*}
\bigskip

Hence $\Delta^{i_1} U^{i_1} > 0$ but $\Delta^{i_2} U^{i_2} < 0$
\bigskip

\alert{Disagreement!} Shareholders do not converge on a common objective
\bigskip

This is generic under incomplete markets: investors cannot trade their differences away
\end{frame}

%====================================================================
\section{Unanimity}
%====================================================================

\begin{frame}\frametitle{Unanimity under complete markets}
\begin{proposition}
Under (C.1)--(C.4), if $\Range \kappa = \R^S$, there exists a unique state-price deflator $\lambda \in \R^S_{++}$ with
\begin{equation*}
\MRS^i = \lambda \quad \text{for every investor } i \in I
\end{equation*}
All investors share identical competitive price perceptions and unanimously agree on
\begin{equation*}
V = \delta_0 + \sum_{s \in S} \lambda_s \delta_s = y_0 + \sum_{s \in S} \lambda_s y_s
\end{equation*}
\end{proposition}
\end{frame}

%--------------------------------------------------------------------
\begin{frame}\frametitle{A by-product: irrelevance of $\beta$}
The firm's value
\begin{equation*}
V = y_0 + \sum_{s \in S} \lambda_s y_s
\end{equation*}
depends on the financial plan $\beta$ only through its market value, which is offset by debt obligations
\bigskip

This is a first version of the \alert{Modigliani--Miller irrelevance theorem}
\bigskip

The financial structure of the firm is \emph{irrelevant} for its value, conditional on the production plan $y$
\end{frame}

%--------------------------------------------------------------------
\begin{frame}\frametitle{Spanning condition}
Complete markets are sufficient but \emph{not necessary} for unanimity
\bigskip

\begin{proposition}[Spanning implies unanimity]
Suppose that for every feasible production change $\Delta y$, there exist $\Delta \theta \in \R^A$ and $\Delta \alpha \in \R$ such that
\begin{equation*}
\Delta y_s = \sum_{a \in A} \kappa_s(a) \Delta \theta(a) + y_s \Delta \alpha,
\quad \text{for all $s \in S$}
\end{equation*}
Then every shareholder agrees on the present-value criterion
\end{proposition}
\end{frame}

%--------------------------------------------------------------------
\begin{frame}\frametitle{Spanning: economic content}
\textbf{Economic content:} any change in the firm's output can be replicated through
\begin{itemize}
\item a combination of asset trades $\Delta \theta$
\item a proportional rescaling $\Delta \alpha$ of existing production
\end{itemize}
\bigskip

Shareholders do not lose hedging opportunities: their valuation collapses to the marketed present value
\bigskip

\emph{References:} Ekern--Wilson (1974), Leland (1974), Magill--Quinzii (2009)
\end{frame}

%--------------------------------------------------------------------
\begin{frame}\frametitle{Strong unanimity: setup}
\textbf{Goal:} Strengthen the result. Under common MRSs, shareholders not only \emph{agree} on the objective; they would not deviate even with full control
\bigskip

\begin{block}{No-default plans}
\begin{equation*}
\mathcal{D} := \bigl\{ (y, \beta) \in Y \times \R^A \colon y_s \geq \kappa_s \cdot \beta, \ \forall s \in S \bigr\}
\end{equation*}
\end{block}

\begin{block}{Counterfactual equity price}
For $\lambda \in \R^S_{++}$ and $(y', \beta') \in \mathcal{D}$:
\begin{equation*}
\widetilde{E}(\lambda, y', \beta') := \sum_{s \in S} \lambda_s [y'_s - \kappa_s \cdot \beta']
\end{equation*}
\end{block}
\end{frame}

%--------------------------------------------------------------------
\begin{frame}\frametitle{Strong unanimity: theorem}
\begin{theoremns}[Strong unanimity]
Under (C.1)--(C.4), at any complete-markets equilibrium with common $\lambda$ and $\MRS^i = \lambda$ for every $i$, suppose the firm has chosen $(y, \beta) \in \mathcal{D}$ to maximize
\begin{equation*}
V(y', \beta') := y'_0 + q \cdot \beta' + \sum_{s} \lambda_s [y'_s - \kappa_s \cdot \beta']
\end{equation*}
Then for every $i \in I$:
\begin{equation*}
(y, \beta) \in \argmax \Bigl\{ U^i(c') \colon (c', \theta', \alpha') \in B^i(q, \lambda, y', \beta'), \ (y', \beta') \in \mathcal{D} \Bigr\}
\end{equation*}
\end{theoremns}
\end{frame}

%--------------------------------------------------------------------
\begin{frame}\frametitle{Strong unanimity: interpretation}
The equilibrium plan is utility-maximizing for every investor over the joint choice of:
\begin{itemize}
\item consumption $c'$
\item portfolio $\theta'$
\item shareholding $\alpha'$
\item firm plan $(y', \beta')$ -- counterfactually controllable
\end{itemize}
\bigskip

No shareholder would prefer to redirect the firm even given full authority over its decisions
\bigskip

This is the strongest possible justification for delegating control to a value-maximizing manager
\end{frame}

%====================================================================
\section{Bibliography}
%====================================================================

\begin{frame}\frametitle{Classical references}
\begin{thebibliography}{8}
\begin{small}
  \beamertemplatearticlebibitems
    \bibitem{Fisher1930}
    Fisher, I.
    \newblock {\em The Theory of Interest}.
    \newblock New York: A. M. Kelley, 1930.
    \bibitem{Dreze1974}
    Dr\`eze, J. H.
    \newblock Investment under private ownership: Optimality, equilibrium and stability
    \newblock In \emph{Allocation Under Uncertainty: Equilibrium and Optimality}, Wiley, 1974.
    \bibitem{EkernWilson1974}
    Ekern, S. and Wilson, R.
    \newblock On the theory of the firm in an economy with incomplete markets.
    \newblock {\em Bell Journal of Economics and Management Science} 5(1), 1974.
    \bibitem{Leland1974}
    Leland, H. E.
    \newblock Production theory and the stock market.
    \newblock {\em Bell Journal of Economics and Management Science} 5(1), 1974.
\end{small}
\end{thebibliography}
\end{frame}

%--------------------------------------------------------------------
\begin{frame}\frametitle{Foundational references}
\begin{thebibliography}{8}
\begin{small}
  \beamertemplatearticlebibitems
    \bibitem{GrossmanHart1979}
    Grossman, S. J. and Hart, O. D.
    \newblock A theory of competitive equilibrium in stock market economies.
    \newblock {\em Econometrica} 47(2), 1979.
    \bibitem{GrossmanStiglitz1980}
    Grossman, S. J. and Stiglitz, J.
    \newblock Stockholder unanimity in making production and financial decisions.
    \newblock {\em Quarterly Journal of Economics} 94(3), 1980.
    \bibitem{Geanakoplos1990}
    Geanakoplos, J., Magill, M., Quinzii, M. and Dr\`eze, J.
    \newblock Generic inefficiency of stock market equilibrium when markets are incomplete.
    \newblock {\em Journal of Mathematical Economics} 19, 1990.
    \bibitem{KelseyMilne1996}
    Kelsey, D. and Milne, F.
    \newblock The existence of equilibrium in incomplete markets and the objective function of the firm.
    \newblock {\em Journal of Mathematical Economics} 25, 1996.
\end{small}
\end{thebibliography}
\end{frame}

%--------------------------------------------------------------------
\begin{frame}\frametitle{Modern developments}
\begin{thebibliography}{8}
\begin{small}
  \beamertemplatearticlebibitems
    \bibitem{TvedeCres2005}
    Tvede, M. and Cr\'es, H.
    \newblock Voting in assemblies of shareholders and incomplete markets.
    \newblock {\em Economic Theory} 26, 2005.
    \bibitem{MagillQuinzii2008}
    Magill, M. and Quinzii, M.
    \newblock Normative properties of stock market equilibrium with moral hazard.
    \newblock {\em Journal of Mathematical Economics} 44, 2008.
    \bibitem{MagillQuinzii2009}
    Magill, M. and Quinzii, M.
    \newblock The probability approach to general equilibrium with production.
    \newblock {\em Economic Theory} 39, 2009.
    \bibitem{CarcelesCoenPirani2009}
    Carceles-Poveda, E. and Coen-Pirani, D.
    \newblock Shareholders' unanimity with incomplete markets.
    \newblock {\em International Economic Review} 50, 2009.
\end{small}
\end{thebibliography}
\end{frame}

%--------------------------------------------------------------------
\begin{frame}\frametitle{Recent contributions}
\begin{thebibliography}{8}
\begin{small}
  \beamertemplatearticlebibitems
    \bibitem{MagillQuinzii2010}
    Magill, M. and Quinzii, M.
    \newblock A co-moment criterion for the choice of risky investment by firms.
    \newblock {\em International Economic Review} 51, 2010.
    \bibitem{BejanBidian2011}
    Bejan, C. and Bidian, F.
    \newblock Ownership structure and efficiency in large economies.
    \newblock {\em Economic Theory} 50, 2012.
    \bibitem{DierkerDierker2011}
    Dierker, E. and Dierker, H.
    \newblock Dr\`eze equilibria and welfare maxima.
    \newblock {\em Economic Theory} 45, 2009.
    \bibitem{BritzHeringsPredtetchinski2013}
    Britz, V., Herings, P.J.J. and Predtetchinski, A.
    \newblock A bargaining theory of the firm.
    \newblock {\em Economic Theory} 54, 2013.
\end{small}
\end{thebibliography}
\end{frame}

%--------------------------------------------------------------------
\begin{frame}\frametitle{Recent contributions (cont.)}
\begin{thebibliography}{8}
\begin{small}
  \beamertemplatearticlebibitems
    \bibitem{BisinGottardiRuta2014}
    Bisin, A., Gottardi, P. and Ruta, G.
    \newblock Equilibrium corporate finance and intermediation.
    \newblock {\em Working Paper}, 2015.
    \bibitem{BisinClementiGottardi2025}
    Bisin, A., Clementi, G. L. and Gottardi, P.
    \newblock Capital structure and hedging demand with incomplete markets.
    \newblock {\em Economic Journal}, 2025.
\end{small}
\end{thebibliography}
\bigskip

\textbf{Textbook treatment of incomplete-markets general equilibrium:}
\bigskip

Magill, M. and Quinzii, M., \emph{Theory of Incomplete Markets}, MIT Press, 1996.
\end{frame}

\end{document}